% Extracto íntegro por residencia del cierre físico-matemático sucesor.
\section{Entrelaceur canonique entre mémoire TPK et appareil}
\label{sec:hmtf2-aparato-memoria}

Soit \(\mathfrak M_n\) l'ensemble fini des registres TPK au niveau \(n\),
comprenant le résultat, la route, la retenue, le feuillet et la mémoire. Son hilbertisation
canonique est
\[
 \mathcal H_{\mathfrak M_n}
 =\ell^2(\mathfrak M_n).
\]
Soit \(\mathcal K_{A,n}\) un espace auxiliaire de l'appareil contenant une
famille orthonormée
\(\{|a_m\rangle:m\in\mathfrak M_n\}\). Nous définissons
\[
 \iota_n:\mathcal H_{\mathfrak M_n}\longrightarrow\mathcal K_{A,n},
 \qquad
 \iota_n|m\rangle=|a_m\rangle.
\]
Pour ne pas perdre de généalogie lorsque la mise à jour visible n'est pas injective,
on utilise le relèvement de registre
\[
 \widehat F_n(m)=(F_n(m),m)\in\mathfrak M_{n+1}.
\]
Il est injectif par construction et définit l'isométrie
\[
 V_n|m\rangle=|\widehat F_n(m)\rangle.
\]
Sur le sous-espace de pointeur, on fixe
\[
 W_n\iota_n=\iota_{n+1}V_n.
\]

\begin{theorem}[Entrelaceur appareil--mémoire TPK]
\label{thm:hmtf2-entrelazador-memoria}
Les applications \(\iota_n\) et \(V_n\) sont des isométries et le carré
\[
 \begin{array}{ccc}
 \mathcal H_{\mathfrak M_n}&\xrightarrow{\ V_n\ }&
 \mathcal H_{\mathfrak M_{n+1}}\\
 {\scriptstyle\iota_n}\downarrow&&\downarrow{\scriptstyle\iota_{n+1}}\\
 \mathcal K_{A,n}&\xrightarrow{\ W_n\ }&
 \mathcal K_{A,n+1}
 \end{array}
\]
commute. L'application \(W_n\) se prolonge en une isométrie de l'appareil et, après
ajout d'un supplémentaire auxiliaire si les dimensions diffèrent, en une application unitaire.
La mise à jour préserve la norme, distingue tous les registres et conserve l'
état précédent dans le registre.
\end{theorem}

\begin{proof}
Les applications envoient des bases orthonormées sur des familles orthonormées. L'identité
du carré se vérifie sur chaque \(|m\rangle\). Toute isométrie finie s'
étend en une application unitaire entre dilatations de même dimension en complétant des
bases orthonormées.
\end{proof}

Les projecteurs de pointeur
\[
 P_m=|a_m\rangle\langle a_m|
\]
définissent l'espérance conditionnelle
\[
 \mathcal D_n(A)=\sum_{m\in\mathfrak M_n}P_m A P_m.
\]
On a
\[
 \mathcal D_n^2=\mathcal D_n,\qquad
 \operatorname{Spec}(\mathcal D_n)\subset\{0,1\}.
\]
La sous-algèbre diagonale des registres est le sous-espace fixe et le supplémentaire
cohérent appartient au noyau. L'écart unitaire de cette projection vaut un.
Sous une conjugaison simultanée de l'appareil, des pointeurs et de l'
entrelaceur, toutes ces identités demeurent exactes. Ainsi est établie
la stabilité finie du registre au sein de la classe covariante.

Lorsque les étiquettes \((y,\alpha)\) de Kraus sont incluses dans
\(\mathfrak M_n\), l'isométrie \(\iota_n\) identifie explicitement la
mémoire TPK à l'environnement de Stinespring, et pas seulement à un espace auxiliaire
abstrait. La construction d'un dispositif matériel, sa décohérence et
son étalonnage constituent une \textsc{reconnaissance extérieure}.


